\section{Cognitive Model}
\label{sec:cognitive-model}

\subsection{From walk objects to cognitive quantities}
\label{sec:cognitive-mapping}

The model treats a graph as a semantic or decision network and the walk as a
deliberation process. \Cref{tab:mapping} lists the correspondence. The mapping
is deliberately minimal: every cognitive quantity is a statistic of the walk,
so the model can be simulated and compared with a classical baseline without
additional free parameters.

\begin{table}[t]
  \centering
  \small
  \begin{tabular}{@{}p{0.34\linewidth}p{0.58\linewidth}@{}}
    \toprule
    Walk object & Cognitive interpretation \\
    \midrule
    Vertex $x\in V$ & concept, evidence state, or option \\
    Edge $(x,y)\in E$ & semantic association or evidence transition \\
    Marked set $M$ & target memory or decision boundary \\
    Initial state $\lvert\psi_0\rangle$ & cue, question, or initial attention \\
    Unitary $W(s)$ or $e^{-iLt}$ & deliberation / evidence accumulation \\
    Measurement & committing to an answer or decision \\
    $p_M(t)$ & target activation or success probability \\
    $T_q=\arg\max_t p_M(t)$ & response time at peak evidence \\
    First threshold crossing & response time under a deadline \\
    $P_i$ for option sets $M_i$ & choice probability \\
    $\HT$, $\HTp$ & expected search difficulty, context difficulty \\
    Spectral gap $\lambda_1$ & deliberation / relaxation speed \\
    Interpolation $s$ & commitment or absorption strength \\
    Measurement rate $p$ & introspection or noise level \\
    \bottomrule
  \end{tabular}
  \caption{Mapping between walk quantities and cognitive quantities. The
  mapping is process-level and mathematical; it is not a claim about the
  physical implementation of cognition.}
  \label{tab:mapping}
\end{table}

\subsection{Modeling assumptions}
\label{sec:cognitive-assumptions}

We make five assumptions explicit.
\begin{enumerate}[leftmargin=1.4em]
  \item \textbf{Process level.} The model describes the statistics of a search
  or decision process. It does not claim that the brain implements a quantum
  walk or a unitary circuit.
  \item \textbf{Interference as multiple paths.} Quantum interference stands
  for the interaction of multiple evidence paths or semantic associations. The
  phase structure is a modeling device, not a physiological statement.
  \item \textbf{Classical baseline.} The baseline is a Markov random walk on
  the same graph with the same initial distribution and the same marked set.
  This isolates the effect of interference from the effect of graph structure.
  \item \textbf{Defined, not fitted, statistics.} Response time, choice
  probability, target activation, entropy, and radius are defined from the walk
  statistics. We do not fit behavioral data in this paper.
  \item \textbf{Graph as prior.} The graph encodes the semantic or decision
  prior. Changing edges or weights changes the predicted context effects.
\end{enumerate}

\subsection{Response time and choice probability}
\label{sec:cognitive-rt}

For a single target, the activation curve is $p_M(t)$. We use two response-time
definitions. The \emph{peak time} is
$T_q=\arg\max_{t\ge1}p_M(t)$, the time at which the evidence for the target is
strongest. The \emph{threshold time} is the first $t$ with
$p_M(t)\ge\theta$, where $\theta$ is a decision threshold. The threshold
definition is closer to a deadline experiment; the peak definition is closer to
an unconstrained search. For decoherence we report the peak success
probability, the success averaged over a uniform deadline, and the probability
that a measurement has already detected the target by the deadline.

For multiple option sets $M_1,\dots,M_k$, the option probabilities at time $t$
are
\begin{equation}
  P_i(t)=\left\lVert(I\otimes\Pi_{M_i})W^t(s)\lvert\psi_0\rangle\right\rVert^2,
  \qquad
  P_i=\frac{P_i(T)}{\sum_jP_j(T)},
\end{equation}
where $T$ is the decision time. Because $P_i(t)$ oscillates, $T$ must come
from a phase-free stopping rule (a random stopping time or a measurement that
detects an option, \cref{sec:methods-decision}); a first-threshold rule makes
the choice depend on the oscillation phase. The choice entropy
$-\sum_iP_i\log P_i$ measures how concentrated the choice distribution is. A
uniform distribution gives maximal entropy; a strong context effect gives low
entropy and a large $\max_iP_i-1/k$ bias.

\subsection{Similarity and context effects}
\label{sec:cognitive-context}

The decision experiments manipulate similarity by coupling two branches at
matching depths with weight $c$. A large $c$ means that the two options are
strongly associated: a walker can move between them without returning to the
hub. In a classical random walk this mainly redistributes probability among
the coupled options; on a symmetric graph the classical choice is unchanged.
In a quantum walk the coupling changes the interference pattern, so the
uncoupled option gains share. This is a
model-level analogue of the similarity effect: the choice
probability of an option depends on the presence and similarity of the other
options, not only on its own value.

The semantic experiments manipulate distance in the same spirit. A cue and a
target in the same community are close in the graph; a cross-community target
requires a longer path. The CTQW and CRW are expected to differ most when the
target is far from the cue, because the quantum walk maintains coherent
amplitude across several paths while the classical walk relaxes locally.

\subsection{Testable predictions}
\label{sec:cognitive-predictions}

The model yields four predictions that can be tested numerically and, in
principle, behaviorally.
\begin{enumerate}[leftmargin=1.4em]
  \item For well-connected graphs and a single target, quantum search time
  scales as $\sqrt{\HT}$ and beats the classical baseline.
  \item Independent targets are easier than clustered targets; a large
  $\HT^+$ penalty requires a structured, not merely clustered, marked set.
  \item Semantic distance changes activation timing; the quantum walk has a
  higher peak for far targets, while the classical walk has a higher early AUC
  for near targets.
  \item Coherence and measurement trade off: the pure walk has the highest
  timed peak, while a partially measured walk that can stop on detection finds
  extended (clustered) targets faster than both the pure and the classical
  walk at moderate deadlines.
\end{enumerate}
These predictions are model-level. Confirming them numerically does not show
that human cognition is quantum; it shows that the quantum-walk model makes
distinct, falsifiable predictions relative to its classical counterpart.
